% =====================================================================
%  Guía práctica: Crea tu primera maquetación de página en ConTeXt
%  Mismo sistema de diseño que las guías anteriores.
%  Compilar con:  context maquetacion.tex
% =====================================================================

% ---------------------------------------------------------------------
% 1. Tipografía y colores
% ---------------------------------------------------------------------
%\definefontfamily [mainface] [rm] [TeX Gyre Pagella]
%\definefontfamily [mainface] [ss] [Latin Modern Sans]
%\definefontfamily [mainface] [tt] [Latin Modern Mono]
%\setupbodyfont     [mainface,11pt]
%\mainlanguage[es]
%
%\definecolor [inkblue]    [r=0.11, g=0.16, b=0.32]
%\definecolor [accent]     [r=0.72, g=0.20, b=0.20]
%\definecolor [accentb]    [r=0.10, g=0.35, b=0.45]
%\definecolor [resultgreen][r=0.13, g=0.42, b=0.24]
%\definecolor [softgray]   [r=0.45, g=0.45, b=0.45]
%\definecolor [panelbg]    [r=0.96, g=0.96, b=0.94]
%\definecolor [resultbg]   [r=0.94, g=0.97, b=0.94]
%\definecolor [codebg]     [r=0.965,g=0.965,b=0.98]
%\definecolor [codeframe]  [r=0.80, g=0.80, b=0.85]
%
%\setupcolors [state=start]
%
%% ---------------------------------------------------------------------
%% 2. Página y márgenes
%% ---------------------------------------------------------------------
%\setuppapersize   [letter]
%\setuplayout      [
%  topspace=2cm, backspace=3cm,
%  header=1cm, footer=1cm,
%  width=middle, height=middle,
%  leftmargin=0cm, rightmargin=0cm,
%]
%\setupwhitespace [medium]
%\setupindenting  [never]
%
%% ---------------------------------------------------------------------
%% 3. Encabezados de sección
%% ---------------------------------------------------------------------
%\definehead [tutorialsection] [section]
%
%\setuphead [tutorialsection]
%  [style=\bf,
%   color=inkblue,
%   number=yes,
%   numbercolor=accent,
%   numberstyle=\bf]
%
%\setuphead [subsection]
%  [style=\bf, color=accentb]
%
%\setuphead [subsubsection]
%  [style=\bf, color=accentb]
%
%% ---------------------------------------------------------------------
%% 4. Bloques de código
%% ---------------------------------------------------------------------
%\definetyping  [contextcode] [
%  option=TEX,
%  space=on,
%  style=\tt,
%]
%
%\setupbackgrounds [contextcode] [background=color, backgroundcolor=codebg]
%
%\setuptyping [contextcode] [
%  frame=off,
%  framecolor=codeframe,
%  rulethickness=0pt,
%  leftmargindistance=1em,
%  rightmargindistance=1em,
%  topspace=0.6em,
%  bottomspace=0.6em,
%]
%
%% ---------------------------------------------------------------------
%% 5. Cajas: Problema, Solución, Explicación, Observación, Resultado
%% ---------------------------------------------------------------------
%\definestartstop [problema] [before={\blank[small]}, after={\blank[small]}]
%\definestartstop [solucion] [before={\blank[small]}, after={\blank[small]}]
%\definestartstop [explicacion] [before={\blank[small]}, after={\blank[small]}]
%\definestartstop [observacion] [before={\blank[small]}, after={\blank[small]}]
%\definestartstop [resultado] [before={\blank[small]}, after={\blank[small]}]
%
%\define[1]\Problema{\blank[small]\startproblema{\bf Problema. }#1\stopproblema\blank[small]}
%\define[1]\Solucion{\blank[small]\startsolucion{\bf Solución. }#1\stopsolucion\blank[small]}
%\define[1]\Explicacion{\blank[small]\startexplicacion{\bf Explicación. }#1\stopexplicacion\blank[small]}
%\define[1]\Observacion{\blank[small]\startobservacion{\bf Observación. }#1\stopobservacion\blank[small]}
%\define[1]\Resultado{\blank[small]\startresultado{\bf\color[resultgreen]{Resultado.} }#1\stopresultado\blank[small]}
%
%% ---------------------------------------------------------------------
%% 6. Interacción
%% ---------------------------------------------------------------------
%\setupinteraction [state=start, color=accentb, contrastcolor=accentb,
%                   title={Maquetación de página en ConTeXt}]
%
%% ---------------------------------------------------------------------
%% 7. Pie de página
%% ---------------------------------------------------------------------
%\setupheadertexts  [][]
%\setupfootertexts  [][{\ss\small\color[softgray]{\pagenumber}}]

\environment env_tuto
% =====================================================================
\starttext

% ------------------------ PORTADA -------------------------------------
\startstandardmakeup[pagestate=start]
  \blank[3*big]
  \blank[medium]
  \blank[medium]
  {\ss\bf\switchtobodyfont[24pt]\color[inkblue]{Tutorial}}\par
  \blank[small]
  {\ss\bf\switchtobodyfont[17pt]\color[inkblue]{Crea tu primera}}\par
  {\ss\bf\switchtobodyfont[17pt]\color[tutorialaccent]{maquetación de página en ConTeXt}}\par
  \blank[medium]
  \blank[small]
  \blank[2*big]
  {\it\color[softgray]{Tamaño, márgenes, encabezados, numeración
   y ajustes con \type{\adaptlayout}.}}
  \vfill
  {\ss\small\color[softgray]{Guía práctica paso a paso}}
\stopstandardmakeup

% ------------------------ RESUMEN ---------------------------------------
\startfrontmatter

\tutorialsection{Resumen}
En este tutorial aprenderemos a controlar la apariencia de las páginas
en ConTeXt, paso a paso, construyendo un documento sencillo pero
completo. Al final, tendrás una página con márgenes, encabezado, pie
de página y numeración personalizados.

\blank[medium]
{\ss\bf\color[inkblue]{Contenido}}
\blank[small]
\blank[small]
\placelist[tutorialsection][criterium=all,alternative=c]

\stopfrontmatter

% ------------------------ CUERPO ----------------------------------------
\startbodymatter

% =====================================================================
\tutorialsection{Prepara tu espacio de trabajo}

Primero, crea una carpeta llamada \tt{mi\_maquetacion} y dentro de
ella, un archivo llamado \tt{documento.tex} con este contenido
inicial:

\blank[small]
\startcontextcod
\starttext

\chapter{Introducción}
Este es el contenido de mi documento. Aquí escribiré el texto principal.

\chapter{Siguiente capítulo}
Más contenido para ver cómo se distribuye en la página.

\stoptext
\stopcontextcod
\blank[small]

Compila con \tt{context documento.tex} y observa el resultado.
Obtendrás un PDF con el formato predeterminado de ConTeXt. Ahora,
vamos a modificarlo.

% =====================================================================
\tutorialsection{Establece el tamaño y orientación de la página}

\Problema{Configurar el documento en tamaño A5, orientación vertical
(retrato).}

\Solucion{Añade esta línea justo antes de \type{\starttext}:}

\blank[small]
\startcontextcod
\setuppapersize[A5, portrait]
\stopcontextcod
\blank[small]

El archivo ahora luce así:

\blank[small]
\startcontextcod
\setuppapersize[A5,portrait]

\starttext
...
\stoptext
\stopcontextcod
\blank[small]

\Resultado{Vuelve a compilar. Notarás que la página es más pequeña
que antes (A5). Has cambiado el tamaño del papel.}

\Explicacion{\type{\setuppapersize} controla el tamaño de la página
final. El primer argumento indica el tamaño lógico (A5) y la
orientación (retrato). Como no hemos dado un segundo argumento, el
tamaño físico es el mismo.}

% =====================================================================
\tutorialsection{Ajusta los márgenes}

\Problema{Queremos que el texto no esté tan pegado a los bordes.
Vamos a crear márgenes más amplios.}

\Solucion{Justo debajo de \type{\setuppapersize}, añade:}

\blank[small]
\startcontextcod
\setuplayout
  [backspace=2.5cm,
   width=fit]
\stopcontextcod
\blank[small]

El preámbulo ahora es:

\blank[small]
\startcontextcod
\setuppapersize[A5, portrait]
\setuplayout[backspace=2.5cm, width=fit]

\starttext
...
\stoptext
\stopcontextcod
\blank[small]

\Resultado{Observa cómo el margen izquierdo (backspace) ha
aumentado. La opción \type{width=fit} hace que el ancho del texto se
calcule automáticamente para que el margen derecho sea igual al
izquierdo.}

\Explicacion{Experimenta: cambia el valor de \type{backspace} a 1cm
o a 4cm. Vuelve a compilar cada vez para ver el efecto. Ajusta el
valor a tu gusto.}

% =====================================================================
\tutorialsection{Añade encabezado y pie de página}

\Problema{Incluir un encabezado con el título del capítulo y un pie
de página con el número de página.}

\Solucion{Primero, añade las opciones para activar el encabezado y el
pie de página en \type{\setuplayout}:}

\blank[small]
\startcontextcod
\setuplayout
  [backspace=2.5cm,
   width=fit,
   header=0.8cm,
   footer=0.8cm,
   headerdistance=0.5cm,
   footerdistance=0.5cm]
\stopcontextcod
\blank[small]

Luego, para personalizar el contenido del encabezado y pie, añade las
siguientes líneas \bf{antes de} \type{\starttext}:

\blank[small]
\startcontextcod
\setupheadertexts[{\getmarking[chapter]}]

\setupfootertexts[{\pagenumber}]
\stopcontextcod
\blank[small]

Ahora, el preámbulo completo es:

\blank[small]
\startcontextcod
\setuppapersize[A5, portrait]
\setuplayout
  [backspace=2.5cm,
   width=fit,
   header=0.8cm,
   footer=0.8cm,
   headerdistance=0.5cm,
   footerdistance=0.5cm]

\setupheadertexts[{\getmarking[chapter]}]
\setupfootertexts[{\pagenumber}]

\starttext
...
\stoptext
\stopcontextcod
\blank[small]

\Resultado{Ahora deberías ver el nombre del capítulo en el encabezado
y el número de página en el pie.}

\Observacion{Comprobación: pasa a la segunda página (con el capítulo
``Siguiente capítulo''). El encabezado debería cambiar
automáticamente al nuevo título. El número de página debería ser
``2''.}

% =====================================================================
\tutorialsection{Personaliza la numeración de páginas}

\Problema{Queremos que el número de página aparezca como ``Página 1''
en lugar de solo ``1''.}

\Solucion{Sustituye la línea de configuración del pie de página por
esta:}

\blank[small]
\startcontextcod
\setupfootertexts[Página \pagenumber]
\stopcontextcod
\blank[small]

\Resultado{Ahora el pie muestra ``Página 1'', ``Página 2'', etc.}

\startexplicacion
\bf Explicación. Si quieres más control, podemos usar
\type{\setuppagenumbering} para cambiar la ubicación y el formato:

\startcontextcod
\setuppagenumbering
  [location={footer,right},
   style={\ss\bf},
   color=darkblue]
\stopcontextcod

Añade esta línea antes de \type{\starttext}. Observa los cambios: el
número se alinea a la derecha en el pie de página, cambia de fuente y
color.
\stopexplicacion

% =====================================================================
\tutorialsection{Distingue páginas pares e impares (para impresión a dos caras)}

\Problema{Configurar el documento para que los márgenes sean
simétricos en páginas pares e impares, como en un libro.}

\Solucion{Modifica \type{\setuppapersize} para habilitar la opción ``a
dos caras'':}

\blank[small]
\startcontextcod
\setuppapersize[A5, portrait]
\setuparranging[2UP]  % Simula un documento a dos caras
\stopcontextcod
\blank[small]

\Observacion{Nota: \type{\setuparranging[2UP]} es solo para
visualización en pantalla. En un documento real impreso, usarías
\type{\setuppagenumbering[alternative=doublesided]}.}

Luego, cambia \type{backspace} por las opciones \type{inner} y
\type{outer} en \type{\setuplayout}:

\blank[small]
\startcontextcod
\setuplayout
  [inner=2.5cm,
   outer=1.5cm,
   width=fit,
   header=0.8cm,
   footer=0.8cm,
   headerdistance=0.5cm,
   footerdistance=0.5cm]
\stopcontextcod
\blank[small]

\Resultado{Ahora, al visualizar las páginas, verás que el margen
interior (cerca del lomo) es más grande que el exterior. En una página
par, el margen ``interior'' estará a la derecha.}

\Observacion{Reflexión: ¿Observas cómo cambia el espacio entre el
texto y el borde de la página al pasar de una página par a una impar?
Esto simula la encuadernación de un libro.}

% =====================================================================
\tutorialsection{Ajustes finos con \type{adaptlayout}}

\Problema{Imagina que una página queda con muy poco texto al final de
un capítulo. Queremos que ocupe un poco más de espacio.}

\Solucion{Al final del documento, justo antes de \type{\stoptext},
añade:}

\blank[small]
\startcontextcod
\adaptlayout[+2]
\stopcontextcod
\blank[small]

\Resultado{Con \type{\adaptlayout}, podemos aumentar o disminuir el
número de líneas en la página actual.}

\Observacion{Esta herramienta es muy delicada y, como es una operación
avanzada, simplemente la dejamos aquí para que sepas que existe.}

\tutorialsection{Ejercicio}

Reproduce y mejora este documento:
% ============================================
% Tamaño de página y configuración básica
% ============================================
\setuppapersize[letter]
\setuppagenumbering[location=,alternative=doublesided]
\setupbodyfont[11pt, palatino]
\mainlanguage[es]
% Desactivar encabezado por defecto; el pie lleva el texto decorativo
\setupheadertexts[]
\setupfootertexts[][{\color[deepblue]{\ss\tfx Hecho con \ConTeXt\ · Página \userpagenumber}}]
% ============================================
% Definición de colores
% ============================================
\setuppapersize[letter]
\definecolor[deepblue] [h=173A5A]
\definecolor[warmgold] [h=C5A66A]
\definecolor[white]    [h=FFFFFF]
\definecolor[deepblue][r=0.10, g=0.20, b=0.40]
\definecolor[warmgold][r=0.85, g=0.65, b=0.30]
\definecolor[softgray][r=0.94, g=0.93, b=0.91]
% ============================================
% Diseño: dejar un margen izquierdo amplio como zona decorativa
% ============================================
\setuplayout[
  backspace=5cm,
  cutspace=2cm,
  width=fit,
  topspace=3cm,
  bottomspace=3cm,
  header=0cm,
  footer=2.0cm,
  height=fit,
  margindistance=1cm,
]
% ------------------------------------------------------------
% CAPA
% ------------------------------------------------------------

\definelayer
  [leftpanel]
  [
    x=0mm,
    y=0mm,
    width=\paperwidth,
    height=\paperheight,
    repeat=yes,
  ]

% ------------------------------------------------------------
% FRANJA AZUL
% ------------------------------------------------------------

\setlayer
  [leftpanel]
  [x=0mm,y=0mm]
  {\framed
    [
      frame=off,
      width=4.6cm,
      height=\paperheight,
      background=color,
      backgroundcolor=deepblue,
      offset=0pt
    ]{}}

% ------------------------------------------------------------
% TÍTULO VERTICAL
% ------------------------------------------------------------

\setlayer
  [leftpanel]
  [x=8mm,y=60mm]
  {\rotate
    [rotation=90]
    {\framed
      [
        frame=off,
        width=105mm,
        height=35mm,
        align=middle,
        offset=0pt
      ]
      {\switchtobodyfont[30pt]
       \bf
       \color[white]{Arte tipográfico}

       \blank[medium]
       \blackrule
         [
           width=30mm,
           height=1pt,
           color=warmgold
         ]

       \blank[small]
       \switchtobodyfont[11pt]
       \tf
       \color[white]{La expresión visual en ConTeXt}
      }}}

% ------------------------------------------------------------
% ACTIVACIÓN DEL FONDO
% ------------------------------------------------------------

\setupbackgrounds
  [page]
  [background=leftpanel]

% ============================================
% Barra decorativa del pie de página
% ============================================
\definelayer[footerbar][
  x=0mm, y=0mm,
  width=\paperwidth,
  height=\paperheight,
  repeat=yes,
]
\setlayer[footerbar][
  x=4.5cm, y=\dimexpr(\paperheight-2cm),
  width=\dimexpr(\paperwidth-6.5cm),
]{
  \blackrule[width=\textwidth,height=2pt,color=warmgold]
}

% ============================================
% Activar las capas de fondo
% ============================================
\setupbackgrounds[page][
  background={leftpanel, footerbar},
  state=repeat,
]

% ------------------------------------------------------------
% DOCUMENTO DE PRUEBA
% ------------------------------------------------------------

%\starttext

\startstandardmakeup
\strut
\stopstandardmakeup

%\definecolumnsetspan[wall][n=3]
%\startcolumnset[n=3, distance=0.33cm]
%\startcolumnsetspan
\definecolumnset[estandard][n=2, balance=yes,distance=0.33cm]
\startcolumnset[estandard]
%\startcolumnsetspan[wide]
\definedfont[Serif at 22pt]
\color[deepblue]{Introducción}
\blank[small]
\definedfont[Serif at 12pt]\em
\color[deepblue]{Maquetar no es solo colocar texto, sino hacer que la página hable por sí misma.}

\blank[big]
\section{¿Para qué a dos columnas?}
El mecanismo de columnas de \ConTeXt\ permite que el contenido fluya de forma natural. La columna izquierda y la derecha se equilibran automáticamente mediante columnset, ideal para una lectura tipo revista.

\blank[medium]
Cada párrafo se asienta sobre una retícula base fija; esta es la promesa subyacente de elegancia en \ConTeXt\: la 
retícula hace que la página respire de manera uniforme.

\blank[medium]
\startframedtext[
  frame=off,
  background=color,
  backgroundcolor=softgray,
  width=\textwidth,
  offset=10pt,
  corner=round,
  framecolor=warmgold,
]
\bf Bloque destacado:
aquí se puede usar \type{\definetextbackground} para lograr un fondo que abarca varios párrafos, ideal para citas o notas complementarias.
\stopframedtext

%blank[medium]
\startsection[title={Símbolos y ornamentos}]
El conjunto de símbolos de \ConTeXt\ nos permite crear detalles visuales refinados. Con\break \type{\definesymbol} se pueden redefinir los símbolos de las listas, aportando un toque de diseño incluso a las enumeraciones.
\stopsection
\stopcolumnsetspan
\stopcolumnset
%\stoptext



\stopbodymatter

% ------------------------ APÉNDICES --------------------------------------
\startappendices

\tutorialsection{¡Has logrado!}

Construir una maquetación de página completa en ConTeXt. Ahora sabes:

\startitemize
\item \bf{Tamaño y orientación}: usar \type{\setuppapersize} para
      definir el formato de la página.
\item \bf{Márgenes}: controlar los espacios en blanco con
      \type{\setuplayout}.
\item \bf{Encabezados y pies}: personalizar su contenido y
      dimensiones.
\item \bf{Numeración}: modificar el estilo, ubicación y formato de
      los números de página.
\item \bf{Páginas pares e impares}: configurar márgenes simétricos
      para impresión a dos caras.
\item \bf{Ajustes locales}: conocer \type{\adaptlayout} para retoques
      finales.
\stopitemize

Este flujo de trabajo es la base para crear documentos con una
apariencia profesional y controlada. A medida que tus proyectos
crezcan, estas herramientas te permitirán adaptar el diseño a
cualquier necesidad.

\stopappendices

% ------------------------ BACKMATTER --------------------------------------
\startbackmatter

\tutorialsection{Próximos pasos}

\startitemize
\item \bf{Experimenta}: prueba a cambiar el encabezado para que
      incluya el número de sección o el título del documento. Usa
      \type{\getmarking[section]} o
      \type{\setupheadertexts[Mi Documento]}.
\item \bf{Explora más opciones}: \type{\setuplayout} tiene muchas más
      configuraciones, como \type{edge} para bordes, \type{margin}
      para márgenes laterales, y \type{lines} para establecer la
      altura del texto en número de líneas.
\item \bf{Diferentes diseños}: aprende a usar
      \type{\definelayout} para tener distintos diseños en diferentes
      partes del documento (por ejemplo, una portada a toda página).
\item \bf{Recuerda}: la maquetación es una herramienta al servicio
      del contenido. Un buen diseño facilita la lectura y la
      comprensión.
\stopitemize

\stopbackmatter

\stoptext